\documentclass[11pt]{article}
\usepackage[margin=0.82in]{geometry}
\usepackage{amsmath,amssymb,amsthm,graphicx,microtype}
\usepackage[hidelinks]{hyperref}
\usepackage[T1]{fontenc}

\newtheorem{theorem}{Theorem}
\newtheorem{lemma}[theorem]{Lemma}
\newtheorem{corollary}[theorem]{Corollary}
\newcommand{\D}{\mathbb D}
\newcommand{\T}{\mathbb T}
\newcommand{\dist}{\operatorname{dist}}

\title{Strip-growth obstructions for rank-one truncated Toeplitz symbols\\
\large A partial answer to Conjectures 3.1--3.2 of arXiv:2107.09115}
\author{Automated mathematical-discovery run \texttt{fa\_banach\_001}}
\date{13 August 2026}

\begin{document}
\maketitle

\noindent\textbf{Status.} Partial result, likely valid, subject to human
review.

\section{Source conjecture and scope}

For an inner function $I$, put $K_I^2=H^2\ominus I H^2$.  If $I$ has an
angular derivative at $\zeta\in\T$, then evaluation at $\zeta$ is bounded on
$K_I^2$, with reproducing kernel
\[
 k_\zeta^I(z)=\frac{1-\overline{I(\zeta)}I(z)}{1-\bar\zeta z}.
\]
The rank-one map $k_\zeta^I\otimes k_\zeta^I$ is a bounded truncated
Toeplitz operator.

On page 10, O'Loughlin asks the following equivalent conjectures
\cite{Source}:
\begin{quote}
Every bounded truncated Toeplitz operator on $K_I^2$ has a bounded symbol if
and only if $I$ is one-component.

Every compact truncated Toeplitz operator on $K_I^2$ has a symbol in
$I C(\T)$ if and only if $I$ is one-component.
\end{quote}

\begin{center}
\includegraphics[width=0.91\textwidth]{figures/open_problem_crop.png}
\end{center}

The forward implication for one-component functions is known
\cite{BBK}.  The converse remains open.  The result below proves a new
endpoint obstruction and applies it to two explicit non-one-component inner
functions.  It does not settle the universal converse.

\section{Proof intuition}

Sarason's criterion for the boundary rank-one operator says that a bounded
symbol would allow one to correct $k_\zeta^I$ by a function $Ih$, $h\in H^2$,
so that the resulting analytic function has bounded real part.  Such a
function maps the disk into a vertical strip and can grow only logarithmically
toward the boundary.

Now approach $\zeta$ very tangentially through points $z_n$ at which
$I(z_n)$ is tiny.  The correction $I(z_n)h(z_n)$ is negligible by the $H^2$
point-evaluation estimate, whereas the kernel is essentially
$(1-\bar\zeta z_n)^{-1}$.  If that reciprocal grows faster than the permitted
logarithm, no bounded symbol can exist.  Zeros make the correction vanish
exactly; tiny values near atoms of a singular measure give a zero-free
version of the same mechanism.

\section{The strip obstruction}

\begin{lemma}[Logarithmic growth in a vertical strip]
\label{lem:strip}
If $F\in H^2$ and $\operatorname{Re}F\in L^\infty(\T)$, then a constant
$C_F<\infty$ exists such that
\[
 |F(z)|\le C_F\log\frac{e}{1-|z|},\qquad z\in\D.
\]
\end{lemma}

\begin{proof}
Let $M=\|\operatorname{Re}F\|_\infty$.  The Poisson representation gives
$|\operatorname{Re}F(z)|\le M$ in $\D$.  If $M=0$, then $F$ is a purely
imaginary constant.  Otherwise put
\[
 G=\frac{F-i\operatorname{Im}F(0)}{M}.
\]
Then $|\operatorname{Re}G|\le1$ and $\operatorname{Im}G(0)=0$.
Carath\'eodory's inequality, applied to either $1+G$ or $1-G$, gives
\[
 |G'(z)|\le \frac{4}{1-|z|^2}.
\]
Integrating on the radius from $0$ to $z=re^{it}$ yields
\[
 |G(z)-G(0)|\le 2\log\frac{1+r}{1-r}.
\]
Since $|G(0)|\le1$, the stated estimate follows after absorbing
$M$ and $|\operatorname{Im}F(0)|$ into $C_F$.
\end{proof}

We use the following individual rank-one criterion of Sarason
\cite[Section 5]{Sarason}; see also \cite[Remark 5.4]{BCFMT}.  If $I$ has an
angular derivative at $\zeta$, then $k_\zeta^I\otimes k_\zeta^I$ has a
bounded symbol only if there is $h\in H^2$ such that
\begin{equation}
 \operatorname{Re}\bigl(k_\zeta^I+Ih\bigr)\in L^\infty(\T).
 \label{eq:sarason}
\end{equation}

\begin{theorem}[Tangential small-value obstruction]
\label{thm:obstruction}
Let $I$ be inner and suppose that $I$ has an angular derivative at
$\zeta\in\T$.  Assume that points $z_n\in\D$ satisfy, with
$d_n=|1-\bar\zeta z_n|$,
\begin{align}
 &z_n\longrightarrow\zeta,\qquad I(z_n)\longrightarrow0,
 \label{eq:small1}\\
 &\frac{|I(z_n)|}{\sqrt{1-|z_n|^2}}=o(d_n^{-1}),
 \label{eq:small2}\\
 &d_n\log\frac{e}{1-|z_n|}\longrightarrow0.
 \label{eq:tangent}
\end{align}
Then the bounded rank-one truncated Toeplitz operator
$k_\zeta^I\otimes k_\zeta^I$ has no bounded symbol.  Consequently it has no
symbol in $I C(\T)$.
\end{theorem}

\begin{proof}
Suppose a bounded symbol exists.  By \eqref{eq:sarason}, choose $h\in H^2$
and set $F=k_\zeta^I+Ih$ with bounded real part.  Lemma \ref{lem:strip}
gives
\begin{equation}
 |F(z_n)|\le C_F\log\frac{e}{1-|z_n|}.
 \label{eq:upper}
\end{equation}
Since $I(z_n)\to0$,
\[
 |k_\zeta^I(z_n)|
 =\frac{|1-\overline{I(\zeta)}I(z_n)|}{d_n}
 \ge \frac{1}{2d_n}
\]
for all large $n$.  The Hardy-space point-evaluation bound and
\eqref{eq:small2} give
\[
 |I(z_n)h(z_n)|
 \le \frac{\|h\|_2|I(z_n)|}{\sqrt{1-|z_n|^2}}
 =o(d_n^{-1}).
\]
Thus $|F(z_n)|\ge(3d_n)^{-1}$ eventually.  This contradicts
\eqref{eq:upper} and \eqref{eq:tangent}.
\end{proof}

\section{A Blaschke example beyond every finite Hardy-power test}

\begin{theorem}[All finite Hardy powers, but no bounded symbol]
\label{thm:blaschke}
For $n\ge2$, let
\[
 \varepsilon_n=e^{-\sqrt n},\qquad
 a_n=(1-\varepsilon_n)e^{i/n},
\]
and let $B$ be the Blaschke product with zeros $(a_n)$.  Then $B$ has an
angular derivative at $1$ and
\[
 k_1^B\in H^p\qquad\text{for every finite }p>1,
\]
but $k_1^B\otimes k_1^B$ has no bounded symbol.
\end{theorem}

\begin{proof}
The series $\sum_n\varepsilon_n$ converges, so $B$ is well defined.  Since
$\varepsilon_n=o(n^{-1})$,
\[
 |1-a_n|\asymp n^{-1},\qquad 1-|a_n|^2\asymp\varepsilon_n.
\]
Hence, for every fixed finite $p>1$,
\begin{equation}
 \sum_{n\ge2}\frac{1-|a_n|^2}{|1-a_n|^p}
 \asymp \sum_{n\ge2} e^{-\sqrt n}n^p<\infty.
 \label{eq:HpB}
\end{equation}
The case $p=2$ is the Ahern--Clark angular-derivative condition, while the
general case and the boundary-kernel criterion give $k_1^B\in H^p$
\cite{AhernClark,Cohn}; see also \cite[Section 2.2]{BCFMT}.

Use $z_n=a_n$ in Theorem \ref{thm:obstruction}.  The first two small-value
conditions are immediate because $B(a_n)=0$, and
\[
 |1-a_n|\log\frac{e}{1-|a_n|}
 \asymp \frac{1+\sqrt n}{n}\longrightarrow0.
\]
The theorem proves the assertion.
\end{proof}

The known fixed-power obstruction says that failure of
$k_\zeta^I\in H^p$ for one $p>2$ rules out a bounded symbol
\cite[Theorem 5.3]{BCFMT}.  Estimate \eqref{eq:HpB} shows that this test is
silent for every finite $p$ in Theorem \ref{thm:blaschke}.

\section{A zero-free atomic singular example}

For a finite positive singular measure $\mu$ on $\T$, write
\[
 S_\mu(z)=\exp\left(-\int_\T\frac{\xi+z}{\xi-z}\,d\mu(\xi)\right).
\]

\begin{theorem}[Zero-free version]
\label{thm:singular}
Let
\[
 \xi_n=e^{i/n},\qquad c_n=e^{-\sqrt n},\qquad
 \mu=\sum_{n\ge2}c_n\delta_{\xi_n},
\]
and put $S=S_\mu$.  Then $S$ is zero-free, has an angular derivative at
$1$, and $k_1^S\in H^p$ for every finite $p>1$.  Nevertheless,
$k_1^S\otimes k_1^S$ has no bounded symbol.
\end{theorem}

\begin{proof}
The measure is finite.  For every finite $p>1$,
\[
 \int_\T\frac{d\mu(\xi)}{|1-\xi|^p}
 =\sum_{n\ge2}\frac{c_n}{|1-\xi_n|^p}
 \asymp\sum_{n\ge2}e^{-\sqrt n}n^p<\infty.
\]
As above, the case $p=2$ gives the angular derivative and the general
Ahern--Clark/Cohn criterion gives $k_1^S\in H^p$.

Set $r_n=1-c_n^2$ and $z_n=r_n\xi_n$.  Because $c_n^2=o(n^{-1})$,
\[
 d_n:=|1-z_n|\asymp n^{-1},\qquad
 \log\frac{e}{1-|z_n|}=1+2\sqrt n.
\]
Thus \eqref{eq:tangent} holds.  The atom at $\xi_n$ gives
\begin{align*}
 |S(z_n)|
 &=\exp\left(-\int_\T
       \frac{1-|z_n|^2}{|\xi-z_n|^2}\,d\mu(\xi)\right)\\
 &\le \exp\left(-c_n\frac{1+r_n}{1-r_n}\right)
 \le \exp(-c_n^{-1}).
\end{align*}
In particular $S(z_n)\to0$, and
\[
 \frac{|S(z_n)|}{\sqrt{1-|z_n|^2}}
 \lesssim \frac{e^{-1/c_n}}{c_n}=o(n)=o(d_n^{-1}).
\]
Theorem \ref{thm:obstruction} completes the proof.
\end{proof}

\section{Limitations, upgrade attempts, and novelty bounds}

Theorems \ref{thm:blaschke} and \ref{thm:singular} provide explicit compact
rank-one witnesses with no $I C(\T)$ symbol and bounded rank-one witnesses
with no bounded symbol.  They therefore verify the conjectured negative
direction for two non-one-component functions.  They do not show that every
non-one-component inner function has the tangential sequence required by
Theorem \ref{thm:obstruction}.

Five focused routes were pursued.  First, the exact Carleson-class condition
$\mathcal C_1(I^2)=\mathcal C_2(I^2)$ was compared with Aleksandrov's
characterization by equality of all embedding classes; no valid exponent
bootstrap was found.  Second, point-mass and Clark-measure tests were used to
seek the boundary derivative criterion, but the basic estimates reduce to
the Clark isometry.  Third, Sarason's rank-one criterion yielded Theorem
\ref{thm:obstruction}.  Fourth, the result was upgraded first to the Blaschke
example and then to the zero-free singular example, both beyond every finite
$H^p$ test.  Fifth, product and model-space direct-sum constructions were
tested as a route to a counterexample to the full conjecture; the needed
$H^1$ endpoint decomposition was not obtained.

A bounded search through 13 August 2026 used arXiv, exact-title and
exact-conjecture searches, the source citations, the phrases ``bounded symbol
one-component'', ``rank-one truncated Toeplitz'', ``weak factorization model
space'', and the later survey/lecture literature.  Recent sources still state
the universal converse as open.  The 2010 paper \cite{BCFMT} proves the
fixed-$p$ necessary condition, but no located source states the strip-growth
criterion or either all-finite-$H^p$ example above.  Novelty is plausible but
provisional pending expert review.

\section{Verification notes}

The only non-elementary inputs are Sarason's individual rank-one criterion
and the Ahern--Clark/Cohn boundary-kernel criterion.  Lemma
\ref{lem:strip} is a direct Carath\'eodory estimate.  In both examples the
series estimates dominate every polynomial because $e^{-\sqrt n}$ does.
The singular estimate uses only the contribution of one positive atom, so
there is no cancellation.  No numerical evidence is used.

\begin{thebibliography}{9}
\bibitem{Source}
R. O'Loughlin,
\emph{Symbols of compact truncated Toeplitz operators},
J. Math. Anal. Appl. 507 (2022), 125819; arXiv:2107.09115.

\bibitem{Sarason}
D. Sarason,
\emph{Algebraic properties of truncated Toeplitz operators},
Oper. Matrices 1 (2007), 491--526.

\bibitem{BCFMT}
A. Baranov, I. Chalendar, E. Fricain, J. Mashreghi, and D. Timotin,
\emph{Bounded symbols and reproducing kernel thesis for truncated Toeplitz
operators}, J. Funct. Anal. 259 (2010), 2673--2701; arXiv:0909.0131.

\bibitem{BBK}
A. Baranov, R. Bessonov, and V. Kapustin,
\emph{Symbols of truncated Toeplitz operators},
J. Funct. Anal. 261 (2011), 3437--3456; arXiv:1009.5123.

\bibitem{AhernClark}
P. Ahern and D. Clark,
\emph{Radial limits and invariant subspaces},
Amer. J. Math. 92 (1970), 332--342.

\bibitem{Cohn}
W. Cohn,
\emph{Radial limits and star invariant subspaces of bounded mean
oscillation}, Amer. J. Math. 108 (1986), 719--749.
\end{thebibliography}

\end{document}
